\section{結論}
\label{sec:conclusion}

本稿では、3D Gaussianシーンから得られるレンダリング教師信号を中核とするコート検出システム
\Methodを提示した。そのフェイルクローズなアライメントは、ホールドアウトで検証したライン観測情報から
メートル尺度のシーン--コート変換を復元する。軌道生成器は、その変換を多様なキャリブレーション済みビューへ
展開する。また、生成されるラベルは、レンダリングごとの人手によるアノテーションなしに、正確な
カメラからコートへの6DoF、内部パラメータ、幾何学的妥当性、レンダラによる支持が確認された画素情報、
および物理キーポイントの同一性を提供する。この正解情報により、単一の大域カメラ仮説を予測し、
微分可能な再投影とcheiralityを介して軽量な密予測ヘッドに結合するモデルを構成できる。

中心となる設計上の教訓は、ここで合成データが最も有用なのは、より多くの画像を模倣するからではなく、
各画像の背後にある潜在幾何を、追加のアノテーションコストを限りなく小さく抑えながら明示するからだという
点にある。局所的な視覚情報が、大域的に整合する単一のコート解釈に寄与するよう学習することで、
本システムは下流の3次元テニス再構成に必要な座標基盤を提供する。本手法を実証的に採用するかどうかは、
整合性損失だけではなく、事前登録された複数シードおよび実データへの転移プロトコルによって
引き続き判断される。
